\documentclass[pdflatex,sn-mathphys-num]{sn-jnl}

\usepackage{graphicx}%
\usepackage{multirow}%
\usepackage{amsmath,amssymb,amsfonts}%
\usepackage{amsthm}%
\usepackage{mathrsfs}%
\usepackage[title]{appendix}%
\usepackage{xcolor}%
\usepackage{textcomp}%
\usepackage{manyfoot}%
\usepackage{booktabs}%
\usepackage{algorithm}%
\usepackage{algorithmicx}%
\usepackage{algpseudocode}%
\usepackage{listings}%
\usepackage{lmodern}%
\usepackage{siunitx}
\usepackage{microtype}
\usepackage{derivative}
\usepackage[version=4]{mhchem}
\usepackage{cleveref}
\usepackage[spanish]{babel}

\renewcommand{\arraystretch}{1.2}
\sisetup{group-digits=true,
  group-separator={\,},
  separate-uncertainty
}

\theoremstyle{thmstyleone}%
\newtheorem{theorem}{Teorema}%
\newtheorem{proposition}[theorem]{Proposición}%

\theoremstyle{thmstyletwo}%
\newtheorem{example}{Ejemplo}%
\newtheorem{remark}{Observación}%

\theoremstyle{thmstylethree}%
\newtheorem{definition}{Definición}%

\raggedbottom
\unnumbered% uncomment this for unnumbered level heads

\begin{document}

\title[Mecanismos de Resonancia Plasmónica]{Mecanismos de Resonancia
Plasmónica y su Aplicación en Biosensado}

\author{\fnm{Julian L.} \sur{Avila-Martinez}}
\email{jlavilam@udistrital.edu.co}

\affil{\orgdiv{Programa Académico de Física}, \orgname{Universidad Distrital
Francisco José de Caldas}, \orgaddress{\city{Bogotá D.C.},
\country{Colombia}}}

\maketitle

\section{Características Fundamentales de la Resonancia de Plasmón}
\label{sec:caracteristicas}

El fenómeno físico de la resonancia de plasmón de superficie exhibe tres
características principales críticas para las arquitecturas de biosensado.
Primero, la interacción entre la radiación electromagnética incidente y la
nanopartícula induce una oscilación colectiva y cuantizada de la nube de
electrones, creando un dipolo oscilante. Este confinamiento genera regiones
localizadas de dimensiones sublongitud de onda donde la amplitud del campo
eléctrico se amplifica en varios órdenes de magnitud. Segundo, la frecuencia
de resonancia depende intrínsecamente de la estructura morfológica de la
nanopartícula y de la permitividad del medio dieléctrico circundante. Esta
dependencia ambiental constituye el mecanismo de transducción para el sensado.
Tercero, el sistema dinámico disipa una fracción de la energía
electromagnética absorbida como calor mediante radiación térmica, una
consecuencia directa de la radiación de frenado durante la oscilación dipolar.

En aplicaciones de biosensado, estas características permiten la determinación
precisa de las afinidades de enlace molecular. Al monitorear los
desplazamientos en la condición de resonancia causados por cambios locales en
el índice de refracción, es posible cuantificar interacciones como la
capacidad de unión entre un rotavirus y proteínas específicas recubiertas con
polietilenglicol. Contrario a un tratamiento estrictamente clásico, el origen
de la resonancia radica en la naturaleza cuantizada del plasmón.

\section{Tipología de las Resonancias Plasmónicas}
\label{sec:tipologia}

Las condiciones de frontera espaciales impuestas al gas de electrones definen
tres regímenes resonantes distintos.

La resonancia de plasmón de superficie localizado involucra excitaciones
confinadas espacialmente a la geometría finita de una nanoestructura. Los
modos no se propagan y se caracterizan por fuertes amplificaciones del campo
local.

Los polaritones de plasmón de superficie representan ondas electromagnéticas
propagantes acopladas a las oscilaciones de densidad de carga en la interfaz
plana entre un dieléctrico y un conductor. Estos estados requieren
condiciones específicas de acoplamiento de fase para ser excitados.

La resonancia de red de superficie emerge cuando las nanoestructuras
individuales se organizan en una base de red periódica. El acoplamiento
difractivo de los plasmones localizados, mediado por anomalías de Rayleigh en
el arreglo periódico, hibrida los modos localizados.

Entre estas modalidades, la resonancia de red de superficie posee la mayor
utilidad para aplicaciones de biosensado. El arreglo periódico suprime el
amortiguamiento radiativo, conduciendo a una respuesta espectral más estrecha
y una sensibilidad mejorada que supera las capacidades de la resonancia
localizada estándar.

\section{Cinemática de la Configuración de Kretschmann}
\label{sec:kretschmann}

La excitación de un polaritón de plasmón de superficie por un campo
electromagnético externo está prohibida en el espacio libre debido a una
violación estricta de la conservación del momento. La relación de dispersión
de un fotón libre es estrictamente lineal, mientras que la curva de dispersión
del polaritón de plasmón de superficie es no lineal. En consecuencia, para
cualquier frecuencia dada, el momento de la onda de superficie excede
estrictamente al del fotón en el espacio libre.

La configuración de Kretschmann elude esta restricción cinemática al alterar
la densidad óptica del medio incidente. Emplea un prisma dieléctrico de alto
índice adyacente a una película metálica delgada. La presencia del dieléctrico
denso incrementa el vector de onda de la luz incidente. Cuando la radiación
incidente impacta la interfaz metálica en un ángulo crítico específico, la
componente paralela del vector de onda del fotón iguala exactamente al vector
de onda del plasmón en la interfaz opuesta. Este logro del acoplamiento de
fase facilita la transferencia resonante de energía del campo óptico al modo
plasmónico propagante.

\end{document}
